\subsection*{MT-M19. Unconditional nuclear convolution and its exact positivity criterion}
Use the canonical quotient, stable block algebras \(J_n\), and
topological isomorphism \(Q\simeq s(J)\) of MT-M18. No GRH,
simplicity, or reciprocal-derivative bound is assumed initially.
For every \(f\in\mathcal E\), convolution \(C_f\) is nuclear on
\(Q\). There is a Hilbert completion \(\mathcal H_B\) obtained
by making an Auerbach basis in each block orthonormal, on which
every \(C_f\) is trace class, with
\[
 \operatorname{Tr}_{\mathcal H_B} C_f
       =\sum_{\rho\in Z}m_\rho\mathcal Mf(\rho).       \tag{M19.1}
\]
The series is absolutely convergent and the trace is independent
of all block and Auerbach choices. Nuclearity is uniform when
\(f\) ranges over a bounded subset of \(\mathcal E\).

Define \(h^\sharp(x)=e^{-x}\overline{h(-x)}\). Then
\[
 \operatorname{Tr}_{\mathcal H_B} C_{h*h^\sharp}\ge0
          \quad\hbox{for every }h\in\mathcal E
 \quad\Longleftrightarrow\quad
 \text{GRH for }L(s,\eta).                              \tag{M19.2}
\]
Under GRH the trace equals
\(\sum_\rho m_\rho|\mathcal Mh(\rho)|^2\).
This concerns the trace; the chosen block Hilbert norm need
not make the operator \(C_{h*h^\sharp}\) positive.

\begin{proof}
Put \(F=\mathcal Mf\) and \(b_n(f)=\sup_{w\in\Gamma_n}|F(w)|\).
These numbers decrease faster than every power of \(2+|n|\),
with each decay bound controlled by finitely many defining
seminorms of \(f\). This follows from the integration-by-parts
estimate in MT-M18 on the one fixed strip containing all
\(\Gamma_n\). Convolution preserves the closed ideal: the
identity \(\mathcal M(f*g)=\mathcal Mf\,\mathcal Mg\) holds
by absolute integrability, and convolution is continuous on
\(\mathcal E\). Explicitly its seminorms are bounded by
\[
 p_{N,j}(f*g)
 \le p_{N,j}(f)\int_{\mathbb R}e^{N|y|}|g(y)|\,dy
 \le 2p_{N,j}(f)p_{N+1,0}(g).
\]
On \(J_n\), the induced map \(B_n(f)\) is multiplication by
the restriction of \(F\). The quotient Banach-algebra norm
therefore gives
\[
                      \|B_n(f)\|\le b_n(f).            \tag{M19.3}
\]

Let \(v_{n,j},\phi_{n,j}\), \(1\le j\le d_n\), be the
Auerbach vectors and dual functionals of MT-M18, all of norm
one. Regard each vector as supported in its own block and
write \(\ell_{n,j}(a)=\phi_{n,j}(a_n)\). The family
\(\ell_{n,j}\) is equicontinuous, since
\(|\ell_{n,j}(a)|\le q_0(a)\). Set
\[
 \lambda_{n,j}=(2+|n|)^{-2},\qquad
 u_{n,j}=(2+|n|)^2 B_n(f)v_{n,j}.
\]
The sum of \(|\lambda_{n,j}|\) is finite because
\(d_n=O(\log(2+|n|))\). For each \(k\),
\[
 q_k(u_{n,j})\le(2+|n|)^{k+2}b_n(f),
\]
so \(u_{n,j}\) is a bounded family in \(s(J)\). The series
\[
 C_f a=\sum_{n,j}\lambda_{n,j}\ell_{n,j}(a)u_{n,j}       \tag{M19.4}
\]
is thus an explicit nuclear decomposition, absolutely convergent
in every seminorm. Its value is the block multiplication map
because the finite Auerbach expansion in each block is exact.
The same estimates are uniform on a bounded set of \(f\)'s,
which proves the asserted uniformity.

Give the scalar Auerbach coordinates their ordinary square-sum
norm, and let \(\mathcal H_B\) be its completion. For each
nonempty block this Euclidean norm lies between
\(d_n^{-1/2}\|a_n\|_n\) and
\(d_n^{1/2}\|a_n\|_n\), since the vectors and coordinate
functionals have norm one. Consequently the inclusion
\(s(J)\to\mathcal H_B\) is continuous and injective, and
has dense image because it contains all finitely supported
block vectors. Every matrix entry of \(B_n(f)\) in these
orthonormal coordinates has modulus at most \(b_n(f)\),
by (M19.3). Expressing it as a sum of rank-one matrix units
therefore gives
\[
       \|B_n(f)\|_1\le d_n^2b_n(f),\qquad
       \sum_n\|B_n(f)\|_1<\infty.                     \tag{M19.5}
\]
The block direct sum is trace class and extends \(C_f\).

Algebraically, polynomial Hermite interpolation identifies
\(J_n\) with
\(\prod_{\rho\in Z\cap\Omega_n}\mathbb C[z]/(z^{m_\rho})\).
In the monomial basis of a local factor, multiplication by
\(F(\rho+z)\) is triangular with its constant term
\(F(\rho)\) on every diagonal entry. Its finite-dimensional
trace is \(m_\rho F(\rho)\). Trace is invariant under a
change of basis, so (M19.5) proves (M19.1). It also proves
the choice-independence stated there. Absolute convergence
follows directly from
\(\sum_n d_n b_n(f)<\infty\). All traces in the statement
are the ordinary Hilbert traces just constructed; no
unjustified trace on an arbitrary locally convex space is used.

Direct substitution gives
\[
 \mathcal Mh^\sharp(s)
       =\overline{\mathcal Mh(1-\bar s)}.              \tag{M19.6}
\]
The unitary Hecke functional equation implies that the zero
involution \(\rho\mapsto\rho^\dagger=1-\bar\rho\)
preserves \(Z\) and its multiplicities. Indeed it relates
\(\Lambda(s,\eta)\) to \(\Lambda(1-s,\bar\eta)\),
while complex conjugation relates the latter character to
\(\eta\). The pole-removal polynomial introduces no zeros
into \(Z\). Formula (M19.6) therefore preserves the full
zero-jet ideal, so \(\sharp\) descends to an antilinear
algebra involution of \(Q\). If GRH holds then \(\rho^\dagger=\rho\),
and (M19.1), (M19.6) give the stated nonnegative sum.

Conversely, suppose \(\rho\ne\rho^\dagger\) is a pair
off the critical line, of common multiplicity \(m\).
They have the same ordinate and hence lie in the same
block. In that block prescribe the constant jet 1 at
\(\rho\), the constant jet \(-1\) at \(\rho^\dagger\),
and the zero jet at every other zero; prescribe zero in
all other blocks. This is a finitely supported element
of \(s(J)\). The right inverse in MT-M18 realizes it as
\(\mathcal Mh\) for some \(h\in\mathcal E\). Formula
(M19.1) then has just two nonzero contributions for
\(h*h^\sharp\), both equal to \(-m\). Its trace is
\(-2m<0\), contradicting the left side of (M19.2).
This proves the equivalence.
\end{proof}

This is a precise zero-side Weil positivity criterion in
the canonical quotient, with an explicit unconditional
nuclear and trace-class realization. It supplies no proof
of positivity independent of the zeros. The trace is the
sum of traces on finite-length spectral algebras; multiple
zeros contribute their lengths. In the programme of
cote~1 and SA23 geo-H-lidskii/geo-H-genus, MT-M21--MT-M22 identify the specified Meyer adelic summation
quotient with this model, and MT-M24 identifies the actual
Masullo smooth quotient. MT-M25 treats the actual line Hilbert
completion and its lost jets. A cohomological realization and
a positivity argument defined independently of the zeros remain separate. Adopted \(A_D\)
supplies (M19.2)'s positive conclusion for Dirichlet
functions; \(M_H\) supplies it conditionally in general.

\subsection*{MT-M20. A weaker positive Hilbert completion and its exact obstruction}
Keep the unconditional canonical quotient of MT-M18.
If all nontrivial zeros are simple, individual Mellin
evaluation induces a continuous injective map with dense
image
\[
       Q\longrightarrow\ell^2(Z),\qquad
       [g]\longmapsto(\mathcal Mg(\rho))_{\rho\in Z}.   \tag{M20.1}
\]
It therefore defines a positive-definite continuous Hilbert
norm on \(Q\). Every convolution extends diagonally to a
trace-class operator on this completion. Under GRH,
\(C_{h^\sharp}=C_h^*\), and \(C_{h*h^\sharp}\) is a
positive operator. No reciprocal-derivative bound is
needed for this weaker Hilbert norm.

More precisely, the following are equivalent:
\begin{enumerate}
\item GRH holds for \(L(s,\eta)\), and every nontrivial
zero is simple.
\item There is a positive-definite continuous inner
product on \(Q\) such that every convolution extends to
a bounded operator on its Hilbert completion and every
\(C_{h*h^\sharp}\) is positive there.
\end{enumerate}
The inner product is required to be faithful on the
\emph{full jet quotient}. Semisimplifying its multiple-zero
local algebras changes this requirement.

\begin{proof}
Point evaluation at any zero in \(\Omega_n\) has norm at
most one on the quotient \(J_n\), by the maximum principle
and the definition of its quotient norm. Thus
\[
 \sum_{\rho\in Z}|\mathcal Mg(\rho)|^2
 \le\sum_n d_n\|(Rg)_n\|_n^2
 \le q_2(Rg)^2\sum_n d_n(2+|n|)^{-4}<\infty.           \tag{M20.2}
\]
This proves continuity. When all zeros are simple, zero
values mean all zero jets, and MT-M18 identifies the kernel
exactly with the closed ideal; hence the map is injective.
Every finitely supported zero-value vector lies in its
image, since it specifies elements of finitely many
\(J_n\)'s and the interpolation right inverse realizes
them. This proves density and (M20.1).

Convolution by \(f\) has diagonal entries
\(\mathcal Mf(\rho)\). Their absolute sum is finite by
the local zero count and rapid strip decay, so its extension
is trace class. Under GRH (M19.6) makes the entries for
\(h^\sharp\) the complex conjugates of those for \(h\).
The adjoint and positivity statements follow. This proves
the implication (1) to (2).

For the converse, finite-support elements of the block
algebras are actual elements of \(Q\), by MT-M18, and
convolution realizes their algebra multiplication. If
\(\rho\ne\rho^\dagger\), let \(e,e^\dagger\) be the
local idempotents equal to 1 in the full local algebras
at those two zeros and zero in every other local algebra.
The involution interchanges them. For \(h=e-e^\dagger\)
one has
\[
                  h*h^\sharp=-(e+e^\dagger).
\]
On the nonzero finite-dimensional invariant subspace
\((e+e^\dagger)Q\), convolution by this element is
\(-\mathrm{id}\). Its quadratic form is strictly negative
on every nonzero vector in any positive-definite inner
product. This contradicts (2), and proves GRH.

Now suppose a critical-line zero has multiplicity
\(m\ge2\). In its local algebra
\(\mathbb C[z]/(z^m)\), the involution is
\(z^\sharp=-z\), by (M19.6). Let \(e\) be that local
idempotent and \(n\) the element represented by \(z\),
both zero in all other local factors. They are realized
by elements of \(Q\). On the finite-dimensional invariant
space \(eQ\), write \(N=C_n\). It is a nonzero nilpotent
operator. For real \(u\), the element \(h_u=e+iun\)
satisfies \(h_u^\sharp=h_u\), and
\[
 C_{h_u*h_u^\sharp}|_{eQ}
                     =I+2iuN-u^2N^2.                \tag{M20.3}
\]
By (2) these operators are positive, in particular
self-adjoint, for both \(u=1\) and \(u=-1\).
Subtracting them shows that \(iN\) is self-adjoint
in the restricted positive-definite inner product.
A self-adjoint nilpotent on a finite-dimensional
inner-product space is zero: diagonalize the Hermitian
matrix, whose real eigenvalues must all vanish.
This contradicts \(N\ne0\). Thus all zeros are simple,
completing the converse.
\end{proof}

The Hilbert norm in (M20.1) can be strictly weaker than
the original Fr\'echet topology. Its existence therefore
does not make point evaluations a Schauder coordinate
system for that topology: MT-M16 still gives the exact
extra leading-coefficient condition for that assertion.
With multiple zeros, evaluation kills the higher-jet
nilpotent directions. A semisimple Hilbert model with
\(m_\rho\) copies of each scalar eigenvalue retains
the multiplicity-weighted trace, but is not a faithful
Hilbert completion of this full jet quotient. No
simplicity consequence of \(A_D\) or \(M_H\) is claimed.
